紧算子 ``不只是'' 连续的这一事实, 促使人们用\emph{完全连续}算子这一术语来指紧算子. 这一术语多年前很常见, 但如今已较少使用.

紧算子与其定义域和值域的维数之间的联系, 比 Lemma~\ref{lem:infinite-dimensional-identity-noncompact} 所暗示的还要紧密.

\begin{Theorem}
\label{thm:finite-dimensional-domain-range-compactness}
设 $(\mathcal{D}(A),A)$ 为从 $X$ 到 $Y$ 的线性算子. 那么有以下结论:\begin{enumerate}\item 若 $(\mathcal{D}(A),A)$ 有界且值域 $\mathcal{R}(A)$ 有限维, 则算子 $(\mathcal{D}(A),A)$ 为紧算子.\item 若定义域 $\mathcal{D}(A)$ 有限维, 则算子 $(\mathcal{D}(A),A)$ 为紧算子.\end{enumerate}
\end{Theorem}

\begingroup
\postdisplaypenalty=10000
\renewcommand{\endProof}{\par\nobreak\hfill$\blacksquare$\par\vspace{0.1cm}}
\begin{Proof}
对于第 1 部分, 给定有界序列 $x_n\in\mathcal{D}(A)$. 由于算子 $(\mathcal{D}(A),A)$ 有界, 序列 $A(x_n)\in\mathcal{R}(A)$ 也有界. 由于 $\mathcal{R}(A)$ 有限维, Bolzano-Weierstrass 定理表明 $A(x_n)$ 有收敛子序列. 因而 $(\mathcal{D}(A),A)$ 为紧算子.

对于第 2 部分, 注意算子值域的维数小于或等于定义域的维数. (若 $\{x_i\}_{i=1}^k$ 为 $\mathcal{D}(A)$ 的一组基, 则 $\{A(x_i)\}_{i=1}^k$ 张成 $\mathcal{R}(A)$.) 而且由 Lemma~\ref{lem:finite-dimensional-domain-operators-bounded}, 任何定义域有限维的算子都有界. 因而利用第 1 部分即可完成证明.
\end{Proof}
\endgroup

\begin{Definition}
\label{def:finite-rank-operator}
若 $A\in\mathcal{L}(X,Y)$ 且 $\mathcal{R}(A)$ 有限维, 则称算子 $A$ 具有\textbf{有限秩}.
\end{Definition}

证明算子紧性的一种常用方法是, 用已知为紧算子的其他算子 (例如有限秩算子) 来逼近它. 使用这种逼近方法时, 通常要用到以下结果.

\begin{Theorem}
\label{thm:norm-limit-of-compact-operators}
设 $A_n\in\mathcal{L}(X,Y)$ 为紧算子序列. 假设 $A_n$ 在算子范数下收敛到算子 $A$. 那么 $A$ 为紧算子.
\end{Theorem}

\begingroup
\postdisplaypenalty=10000
\renewcommand{\endProof}{\par\nobreak\hfill$\blacksquare$\par\vspace{0.1cm}}
\begin{Proof}
使用 ``对角线序列'' 论证. 给定有界序列 $\{x_n\}_{n=1}^\infty\subset X$. 由于 $A_1$ 为紧算子, 序列 $A_1(x_n)$ 有收敛子序列. 将这一子序列记为 $\{A_1(x_{1,n})\}_{n=1}^\infty$. 现在, 由于 $\{x_{1,n}\}_{n=1}^\infty$ 有界且 $A_2$ 为紧算子, 可见 $A_2(x_{1,n})$ 有收敛子序列. 将其记为 $\{A_2(x_{2,n})\}_{n=1}^\infty$. 重复这一过程, 从子序列中进一步抽取子序列, 使得当 $j<k$ 时, $\{x_{k,n}\}_{n=1}^\infty$ 为 $\{x_{j,n}\}_{n=1}^\infty$ 的子序列, 且 $\{A_k(x_{k,n})\}_{n=1}^\infty$ 收敛. (回忆, 由于 $\{A_k(x_{k,n})\}_{n=1}^\infty$ 收敛, 它是 Cauchy 序列.)

现在考察对角线序列 $\{x_{n,n}\}_{n=1}^\infty$, 记 $z_n:=x_{n,n}$. 注意这确实是原序列 $x_n$ 的子序列. 断言 $A(z_n)$ 是 Cauchy 序列, 因而收敛. (由于 $x_n$ 为任意有界序列, 这就完成了证明.) 给定 $\varepsilon>0$. 注意对于任意 $i$, $j$ 和 $k$, 有
\begin{equation}\label{eq:compact-operator-limit-three-term-estimate}
\begin{aligned}
\|A(z_i)-A(z_j)\|&\leq\|A(z_i)-A_k(z_i)\|+\|A_k(z_i)-A_k(z_j)\|\\
&\quad+\|A(z_j)-A_k(z_j)\|.
\end{aligned}
\end{equation}
由于 $z_n$ 为有界序列, 且 $A_k\to A$ 于算子范数下成立, 可取充分大的 $k$, 使得对于序列 $z_n$ 的每个元素均有
\begin{equation}\label{eq:compact-operator-limit-uniform-tail-estimate}
\|A(z_n)-A_k(z_n)\|<\varepsilon/3
\end{equation}
现在注意, 对固定的 $k$, 序列 $A_k(z_n)$ 是 Cauchy 序列. 这是因为 $\{z_n\}_{n=k}^\infty$ 是 $\{x_{k,m}\}_{m=1}^\infty$ 的子序列. 因而可取充分大的 $i$ 和 $j$, 使得
\begin{equation}\label{eq:compact-operator-limit-subsequence-cauchy-estimate}
\|A_k(z_i)-A_k(z_j)\|<\varepsilon/3.
\end{equation}
将 \eqref{eq:compact-operator-limit-three-term-estimate} 与 \eqref{eq:compact-operator-limit-uniform-tail-estimate} 和 \eqref{eq:compact-operator-limit-subsequence-cauchy-estimate} 结合, 即完成证明.
\end{Proof}
\endgroup
